Dear Members of the RegLab Search Committee,

I am submitting my application for the RegLab Postdoctoral Fellowship. 
I am a final-year PhD candidate in Computer Science at Northwestern University, advised by Jessica Hullman, with expected completion in 2027. 
% My research asks how AI can measurably improve the decisions of human experts. 
I combine statistical decision theory, machine learning, and behavioral experiments to realize \emph{human-AI complementarity} and measurably improve experts' decisions.
% My research could offer a compelling extension on the RegLab's vision on measuring AI's societal consequences and developing transparent AI assistance for consequential decisions.
This agenda aligns with RegLab's focus on evaluating AI's societal consequences and developing AI assistance for consequential decisions.

I use decision theory to evaluate human-AI team's behavior and design algorithms for improving human-AI decisions.
% It is highly aligned with the goals of RegLab on AI's societal consequences and transparent AI-assistance system in high-stake decision-making.
At FAccT 2024, I developed a Bayesian benchmark that decomposes the losses of human-AI teams into losses from following AI too often or too rarely and losses from failing to identify the cases where its advice is most valuable.
My ICLR 2026 and ICML 2026 work evaluates models and explanations by the additional decision value they provide given what another agent already knows.
For the explanation evaluation, we separate the benefits of explanations by the information that can be used independently in explanations and additional information available in the underlying task information.
Building on these foundations, I developed ComplLLM, accepted at NeurIPS 2026, to post-train LLMs to extract interpretable signals that complement another agent's decisions. 
We demonstrate this framework on real-world decision tasks.
Domain experts judged these signals useful, and a crowdsourced study showed that they improved participants' decision accuracy.
% and show that these complementary signals extracted by LLMs are both perceived as useful by domain experts and beneficial to crowd-sourced human decision-makers.

I am particularly interested in RegLab's interdisciplinary research and its vision on bring AI techniques to the public sector.
The distinction proposed by the new paper of Spangher et al. between verifiability and articulability in predicting creative preferences suggests a question for my research: which aspects of expert judgment can AI assistance convey explicitly, and which require other forms of collaboration?
% I just came across the paper of \href{https://arxiv.org/pdf/2610.03025}{Verifiable, Articulable, and Tacit Components of Preference} by Spangher et al. on arXiv and was really attracted by this idea of spliting the gaps of verifiability and articulability in predicting creative preferences.
% This shares a question in my research: which aspects of expert judgment can AI assistance convey explicitly, and which require other forms of collaboration?
It shares some parallel ideas in my FAccT 2024 paper on reliance and ICML 2026 paper on explanation, where we also focused on decomposing the gaps by more interpretable components, and targeted on giving consequential suggestions by understanding the sources of errors better.


At RegLab, I would extend this agenda to human-AI collaboration in creative and open-ended tasks.
I would investigate whether AI can improve expert creators' work by providing information that complements their judgments.
% Using expert-only performance as a baseline, I would test how these benefits vary with the articulability and verifiability of task criteria.
This could provide guidance on when and how AI assistance improves human-AI collaboration.
% In my opinion, creativity tasks (or open-ended tasks) should have large potential of human-AI complementarity, where AI by design does not have a high capability on creativity while high-performance human creators can achieve high performance (such as the tacit knowledge discussed in Spangher et al.).
% There could be a lot of interesting works on how to exploit this potential and better guidances on providing AI assistance to human creators.

I appreciate your consideration and look forward to hearing from you.

Sincerely,\\
Ziyang Guo
